\SourceLeaf{291}%
\[
\begin{aligned}
\sin(Q+q)&=\sin Q\cos\tfrac{1}{2}(h-\delta)+\cos Q\sin\tfrac{1}{2}(h-\delta)\cos\chi\\
\sin(Q-q)&=\sin Q\cos\tfrac{1}{2}(h+\delta)-\cos Q\sin\tfrac{1}{2}(h+\delta)\cos\chi
\end{aligned}
\]
and their sum divided by \(2\cos Q\) gives
\[
\operatorname{tang}Q\mathord{.}(\cos q-\cos\tfrac{1}{2}h\mathord{.}\cos\tfrac{1}{2}\delta)=-\cos\tfrac{1}{2}h\sin\tfrac{1}{2}\delta\cos\chi
\]
\indent Since \(\cos q-\cos\tfrac{1}{2}h\mathord{.}\cos\tfrac{1}{2}\delta\) is evidently a quantity of second order, \(\sin\tfrac{1}{2}\delta\cos\chi\) and \(\delta\cos\chi\) will also be of that order. Hence, since \(\lambda\) contains the factor \(\cos\chi^2\), \(\lambda hh\delta\) is of fourth order and \(\lambda h\delta\delta\) of fifth order. This justifies dropping these terms.

Under the stated assumption, the final result of this expansion is therefore given by the following formulas, with the arcs themselves written in place of the tangents of \(\psi^0,\psi'\)
\[
\begin{aligned}
\psi^0&=-\frac{ee\cos P\sin P\sin\chi\cos\chi^2h^3}{12\cos\varphi\cos\theta}\\
\psi'&=+\frac{ee\cos P\sin P\sin\chi\cos\chi^2h^3}{12\cos\varphi\cos\theta}
\end{aligned}
\]
\Article{16}
The calculation of the triangulation network on the sphere has three main parts:
\begin{list}{}{\setlength{\leftmargin}{1.8em}\setlength{\labelwidth}{1.4em}\setlength{\labelsep}{.4em}\setlength{\itemsep}{0pt}}
\item[1)] adjusting the angles according to all the condition equations arising from the structure of the network.
\item[2)] calculating all the sides of the triangles.
\item[3)] \begingroup\rightskip=0pt plus 1em
determining the longitudes and latitudes of the triangulation points, together with the orientations of the triangle sides that leave each point.\par\endgroup
\end{list}
\hspace*{\parindent}The longitudes and latitudes on the sphere are then converted to the true longitudes and latitudes on the spheroid. The longitudes are divided by the constant divisor \(\alpha\), and the latitudes are found using the auxiliary table attached here, or another table specially constructed in a similar way if there is reason to choose a different standard parallel.

I pass over the first two operations, which rest on familiar principles, and add a few remarks here about the third. It comes down to solving the following problem\footnote{Since this problem is treated here on its own, some letters can be used here with meanings different from those given above without causing any difficulty.}: from the length, expressed in arc units, of
